

\documentclass[sigconf]{acmart}
%%
%% \BibTeX command to typeset BibTeX logo in the docs
\AtBeginDocument{%
  \providecommand\BibTeX{{%
    Bib\TeX}}}

\setcopyright{acmlicensed}
\copyrightyear{2026}
\acmYear{2026}
\acmDOI{XXXXXXX.XXXXXXX} 

\acmConference[MUM '26]{The 25th International Conference on Mobile and Ubiquitous Multimedia (MUM 2026)}{December 2026}{Glasgow, Scotland, United Kingdom}
\acmISBN{978-1-4503-XXXX-X/26/04}

\usepackage{xcolor}
\input{cc-latex}

\begin{document}
\title{Understanding the Impact of Multi-Agent Personas on Stigmatised Disclosures}

\author{Truc Nguyen}
\email{truc.t.nguyen@oulu.fi}
\affiliation{%
  \institution{University of Oulu}
  \city{Oulu}
  \country{Finland}
}

\author{Aku Visuri}
  \email{aku.visuri@oulu.fi}
\affiliation{%
  \institution{University of Oulu}
  \city{Oulu}
  \country{Finland}
}

\author{Dániel Szabó}
  \email{
daniel.szabo@oulu.fi}
\affiliation{%
  \institution{University of Oulu}
  \city{Oulu}
  \country{Finland}
}

\author{Simo Hosio}
\email{simo.hosio@oulu.fi}
\affiliation{%
  \institution{University of Oulu}
  \city{Oulu}
  \country{Finland}
}




% %% This command updates the header on even-numbered pages
\renewcommand{\shortauthors}{Nguyen, et al.}

\begin{abstract}
When individuals face mental health struggles, the fear of judgment often leads them to mask their reality, using guarded language to protect their identity. While conversational agents (CAs) are often deployed to provide anonymity, the structural dynamics of the interaction, specifically whether the exchange is dyadic or communal, play a critical role in determining user trust. This structure can feel evaluative, making users less likely to be honest. In this paper, we investigate whether a multi-agent system can lower the barrier to disclosure by shifting the interaction from a singular focus on the user to a communal social context. We use a mixed-factorial design to test how users respond to discussions with multiple peer agents compared to a standard single-agent interaction.

\end{abstract}

\begin{CCSXML}
<ccs2012>
   <concept>
 <concept_id>10003120.10003121.10003122.10003334</concept_id>
       <concept_desc>Human-centered computing~User studies</concept_desc>
       <concept_significance>500</concept_significance>
       </concept>
   <concept>
       <concept_id>10003120.10003121.10003124.10010870</concept_id>
       <concept_desc>Human-centered computing~Natural language interfaces</concept_desc>
       <concept_significance>300</concept_significance>
       </concept>
   <concept>
       <concept_id>10010405.10010444.10010449</concept_id>
       <concept_desc>Applied computing~Health informatics</concept_desc>
       <concept_significance>300</concept_significance>
       </concept>
 </ccs2012>
\end{CCSXML}

\ccsdesc[500]{Human-centered computing~User studies}
\ccsdesc[300]{Human-centered computing~Natural language interfaces}
\ccsdesc[300]{Applied computing~Health informatics}

\keywords{conversational agents, mental health, stigma markers, self-disclosure, user study}


\maketitle

\section{Introduction}
Disclosing a mental health struggle is not always easy. In everyday life, people rarely share their deepest struggles immediately, but instead, they ``test the waters'' by using vague language, hedging their statements, and waiting for signs of safety before saying more \cite{Ammari2026RedditInfertility, Saha2019LGBTQ}. 
We developed this mechanism naturally in order to avoid being socially stigmatised. Recently, conversational agents (CAs), due to their excellent ability to provide anonymity and carry structured interactions, have been explored as a method to deal with fear of stigmatisation and support comfortable disclosure. 

Further, whether the exchange is dyadic or communal plays a critical role in determining user trust. Conversational systems with a single-persona can inadvertently resemble a clinical interview rather than a shared confidence \cite{Lee2023SocialContact, Song2025Cooperation}, limiting comfortable self-disclosure. Most current mental health agents, for example, are limited to one-on-one interaction with the users. In such exchanges, the user is the sole focus of the agent's inquiries. This interaction may trigger evaluative anxiety, causing users to become even more guarded and socially performative in their responses \cite{Cui2024ChatbotStigma}. When an interaction feels like a formal assessment, users often revert to ``masking'', filtering their words to protect themselves even when they know their interlocutor is an AI \cite{Meng2026CHALET}.

In this study, we design and evaluate an intervention grounded in multi-agent social modelling to overcome this barrier. Prior literature suggests that when a user observes a group of peer personas discussing their own struggles first, the perceived social risk of disclosure should drop \cite{Sien2025DigitalHumanLibrary}. We are testing whether allowing the user to witness a simulated discussion between peer bots themselves first can provide a clearer, safer social context for participation.

To evaluate this, we are conducting a  study using a controlled mixed-factorial design where we compare the number of peer agents in the discussion and various stages of the conversation when the participant is invited to participate and measure how these conditions affect the user's subsequent disclosure. We are measuring the quality of the interaction through ``stigma markers.'' These markers, drawn from recent computational work on latent themes of stigma, allow us to measure whether social modelling actually helps a user drop their guard. Ultimately, we aim to demonstrate the benefits and some of the design considerations of shifting away from purely dyadic dynamics toward multi-agent peer personas for building systems that people can truly trust with their most sensitive disclosures.

In this paper, we will answer three Research Questions (RQs) designed to explore the effects of multi-agent social modelling:

\begin{itemize}
    \item \textbf{RQ1: How does peer disclosure by scripted agents shape what participants disclose about professional and academic struggles?} We tackled this research question through thematic analysis of conversation transcripts with three conversation setups. Through analysis of the qualities, such as the depth of the participant's self-disclosures the three conditions, and tracking whether participants' behaviour are influenced by the Social Modelling approach, we can determine the extent to which our social modelling approach really helped participants safely self-disclose.
    
    \item \textbf{RQ2: How does participant disclosure unfold across the three escalating conversational stages?} This RQ is approached by performing a stage-by-stage qualitative analysis of the three responses to (using the Baseline condition as the baseline trajectory) alongside the quantitative analysis of the five-point comfort .
    
    \item \textbf{RQ3: How do participants experience being invited to disclose alongside disclosing peers?} We will measure this by comparing post-chat BFNE scores \cite{Leary1983BFNE} across conditions and analyzing post-interaction reflection items, where a finding that ``it made no difference'' is treated as a valuable insight rather than a failure.
\end{itemize}


\section{Related Work}

\subsection{The Role of Stigma in Mental Wellbeing}
Mental health stigma operates as a profound systemic and cultural barrier to effective care, significantly limiting professional support and help-seeking behaviours \cite{Sehgal2025AdolescentChatbots}. This pressure is multifaceted; individuals often face external social stigma, which can become internalised as self-stigma or minority stress \cite{Saha2019LGBTQ, Sehgal2025AdolescentChatbots}. 

In digital communities, users frequently manifest these barriers through an avoidance of emotional expression, driven by traditional ideologies or the fear of judgment \cite{Ammari2026RedditInfertility}. Consequently, users may mask their true mental states to protect their social identity. To effectively foster help-seeking behaviour, supportive technologies must be designed to overcome these initial reservations and lower the barriers of shame \cite{Sien2025DigitalHumanLibrary}.

\subsection{Stigma Markers in Computational Interaction}
In computational interaction, understanding user mental states requires analysing user-generated text for linguistic markers that correlate with psychological distress \cite{Rissola2021MentalStateSurvey}. Individuals experiencing high levels of stigma rarely speak plainly; instead, they filter their language, producing specific psycholinguistic markers, such as words related to hopelessness or self-blame, that indicate their underlying struggles \cite{Saduakassova2026SuicidePrevention}. 

For instance, Meng et al. successfully operationalised these markers by uncovering latent themes of stigma across cognitive, emotional, and behavioural dimensions \cite{Meng2026CHALET}. Analysing the language of marginalised or highly stressed groups on social platforms has demonstrated that these subtle cues reveal the presence of perceived and internalised stigma \cite{Saha2019LGBTQ}. Furthermore, analyses of human interactions with AI companions highlight how users express emotional trajectories involving loneliness and depression through language \cite{Yuan2026Replika}. By focusing on these established cues, we can computationally evaluate if our multi-agent design effectively encourages users to stop masking their feelings.

\subsection{Conversational Agents and Stigma Reduction}
A significant challenge in conversational agent design is managing the perceived evaluative tone of the system \cite{Cui2024ChatbotStigma}. In mental health contexts, the effectiveness of these agents is heavily dictated by the power dynamic of the interaction. 
While chatbot-based social contact has shown promise in reducing the desire for social distancing \cite{Lee2023SocialContact}, standard dyadic interactions often default to a ``vertical'' exchange, resembling a one-way clinical interview rather than a shared conversation. 
As \citeauthor{Cui2024ChatbotStigma} \cite{Cui2024ChatbotStigma} highlight, when a bot assumes a highly interpretative or ``expert'' persona, it can inadvertently undermine its own credibility and prompt users to become more guarded in their responses. 

\subsection{Multi-Agent Systems and Social Modelling}
To move beyond the limitations of single-agent interactions, we propose a design rooted in group-based social modelling. Instead of forcing a direct, high-pressure interrogation, we leverage the dynamics of Socially Interactive Agents (SIAs) to create a shared, low-stakes social space \cite{Lugrin2021HandbookVol1, Lugrin2022HandbookVol2}. By allowing a user to observe a simulated dialogue among multiple peer bots, we provide a ``social script'' that effectively lowers the perceived stakes of the interaction. This approach addresses the ``initial reservations'' users often feel when approaching mental health tools, transforming the user from a solitary subject into an observer of a safe social norm.

This ``gentle introduction'' to disclosure is supported by the concept of digital storytelling libraries, where witnessing the experiences of others can significantly boost a user's self-efficacy and willingness to seek help \cite{Sien2025DigitalHumanLibrary}. This design choice mirrors the organic behaviour found in real-world digital support networks, where experienced ``navigators'' help newcomers overcome the silence of self-stigma by modelling how to share and process trauma \cite{Ammari2026RedditInfertility}. By modelling vulnerability through a group of distinct personas rather than a single authority figure, we aim to offer users the ``social permission'' they need to drop their identity curation and engage in more authentic, unfiltered conversation.

\section{Research Methods}
Our study employed a mixed-factorial experiment to test how three different social configurations of conversational agents influence the depth and quality of disclosures which are normally perceived as stigmatised. To achieve this, we employed a mixed-methods experimental design (on what population) to track how self-disclosure unfolds over three stages of the conversation.

\subsection{Experimental Design}

Using a $3 \times 3$ mixed-factorial design, we isolated the effects of our two primary independent variables: social configuration (between-subjects) and conversational stage (within-subjects).

\begin{itemize}
    \item \textbf{Social Configuration (Between-Subjects):} Participants are randomly assigned to one of three agent configurations. To evaluate whether those observing a communal self-disclosure dialogue experience less fear of stigmatisation and judgement than  those partaking in a standard dyadic exchange, we test three configurations of peer agents:
    \begin{itemize}
        \item \textbf{Baseline (no peers):} A pure dyadic interaction with only the host facilitator (Vieno), featuring no peer disclosure.
        \item \textbf{Single Peer (one peer):} The host facilitator joined by one disclosing peer agent (Sam).
        \item \textbf{Multi-Peer (two peers):} The host facilitator joined by two disclosing peer agents (Sam and Charlie).
    \end{itemize}
\end{itemize}

\noindent\textbf{Conversational Stage (Within-Subjects):} To measure how disclosure evolves over time, we observe cognitive masking and real-time comfort as participants progress through three distinct conversational stages: \textit{Early}, \textit{Mid}, and \textit{Late}, which we describe in \autoref{subsubsection:conv_structure}. 

\subsection{The Interaction Flow: Hybrid Social Modelling}

To evaluate the effect of social modelling on self-disclosure, our system employs a hybrid conversational architecture. Rather than using fully generative multi-agent systems that introduce uncontrolled variables, the chat environment consists of two distinct agent roles: a central facilitator (Vieno) and up to two peer agents (Sam and Charlie).

\subsubsection{The Chatbots}

Vieno is driven by a live Large Language Model (specifically Anthropic's Claude Sonnet 5 via OpenRouter) configured with a system prompt (see \autoref{appendix:systemprompt}) to act solely as a neutral, proactive host. To prevent the agent from acting as an unauthorized therapist or breaking the peer-support illusion, Vieno is programmatically constrained to brief, non-interpretative acknowledgments, forbidding it from offering advice, evaluation, or introducing new topics. If generation times out or encounters errors, the system falls back to a neutral line (``Thanks for sharing that.''), ensuring the participant's session is never interrupted.

In contrast to Vieno, the peer agents (Sam and Charlie) are entirely scripted deterministically to deliver consistent, targeted disclosures. To preserve the illusion of a natural, synchronous conversation, both the live LLM host, Vieno and the scripted peers incorporate a word-count-based typing delay (averaging 5 to 7 seconds) before sending messages.

\subsubsection{The Conversation Structure}
\label{subsubsection:conv_structure}

Drawing from the escalating intimacy model of the Fast Friends protocol \cite{Aron1997FastFriends}, the interaction is structured across three distinct stages:

\textbf{1) Early conversation} Following a brief introduction, Vieno asks a formal, low-stakes question (e.g., ``What are the typical challenges you face in your daily professional or academic life?''). In the multi-agent conditions, the peer bots provide superficial answers (e.g., juggling tasks) before the user gets a chance to share. In the baseline condition, the participant answers the question directly without any preceding peer buffer.

\textbf{2) Mid-conversation} Vieno steers the conversation toward moderately stigmatised experiences, asking about a recent moment of unspoken overwhelm. In the two conditions including peer bots, the peer bots model vulnerability by disclosing feelings of burnout and fear of professional judgment. After observing this peer behavior, the user is invited to share again.

\textbf{3) Late conversation} The conversation progresses to even deeper emotional reflections. Vieno asks a more personal question that is presented to the participant as optional (e.g. ``Is there a specific professional or academic challenge you'd feel comfortable going into?'') By this stage, the user in peer chatbot conditions have had maximum exposure to the social norms established by the peer agents, allowing us to measure whether they choose to engage in deep disclosure or remain guarded.

\subsection{Measures (Dependent Variables)}

Considering the scope of this study, we evaluate the intervention focusing on thematic and Linguistic Inquiry and Word Count (LIWC) analysis to capture the nuances of user disclosure, paired with psychometric scales to track shifts in psychological states across the three conversation stages. The complete breakdown of measures are as follows:

\begin{itemize}
    \item \textbf{Thematic Coding of Stigma Markers:} Building on the dimensions identified by \citet{Meng2026CHALET}, we defined stigma markers as specific linguistic choices or themes that suggest the presence of internalised shame and fear of judgment. We conducted reflexive thematic analysis \cite{Braun2006Thematic} involving multiple coders to identify the specific conversational stage at which users lower their reservations. Intercoder reliability was evaluated using Cohen's Kappa \cite{Cohen1960Kappa}.
    
    \item \textbf{Cognitive Masking (using LIWC):} To quantitatively evaluate cognitive masking and emotional expression, we analysed the chat transcripts using LIWC (version 1.12.0). Specifically, we targeted well-established dimensions such as first-person singular pronouns (e.g., ``I'', ``me'', which correlate with genuine emotional disclosure) versus cognitive processing words and third-person pronouns (which indicate distancing). See the complete breakdown of LIWC variables measured in \autoref{tab:liwc_variables}. These linguistic markers allow us to track masking behaviors and emotional tone across the three conversational stages, supplementing basic metrics like word count.
    
    \item \textbf{Conversational Comfort:} To capture the immediate psychological impact of the social modelling on the comfort of the participant, we administer a 1-item self-reported comfort measure immediately following each conversational stage (Early, Mid, and Late). While single-item measures inherently lack the internal reliability of full inventories, this design is a deliberate and necessary compromise for ecological validity. Deploying a lengthy questionnaire mid-chat would severely disrupt the natural conversational flow and artificially spike user anxiety. This approach provides an unobtrusive temporal map of the user's psychological safety in the moment.
    
    \item \textbf{Evaluative Anxiety:} This construct is defined as a user's apprehension about being evaluated negatively by others. We assess this using a state-modified Brief Fear of Negative Evaluation Scale (BFNE) \cite{Leary1983BFNE}, administered as a post-interaction summary measure.
    
    \item \textbf{Baseline Covariates:} To ensure that our dependent variables are measuring the effect of the multi-agent intervention rather than pre-existing trait differences, participants complete a pre-interaction survey. This includes the Distress Disclosure Index (DDI) \cite{Kahn2001DDI} to control for baseline personality variances, an adapted 5-item version of the Stigma Scale for Receiving Psychological Help (SSRPH) \cite{Pinto2014SSRPH}, and the ``Risks'' factor of the Artificial Intelligence Attitude Scale (AIAS) \cite{Aktay2024AIAS} to isolate AI-specific skepticism.
\end{itemize}


\begin{table}
    \centering
    \begin{tabular}{ccc}
         &  & \\
         &  & \\
         &  & \\
    \end{tabular}
    \caption{Caption}
    \label{tab:liwc_variables}
\end{table}

\subsection{Participant Recruitment and Context}
To evaluate the social-modelling based approach with sufficient statistical power, we aimed to recruit 180 participants. A priori power analysis was conducted using G*Power \cite{Faul2007GPower}. To detect a medium effect size ($f=0.25$) in a mixed analysis of variance (Mixed ANOVA) with 80 percentage power at an alpha level of $0.05$, the calculation indicated a required minimum total sample size of 42 participants. By expanding our recruitment to approximately 180 participants (randomly assigned to yield roughly 60 participants per Baseline, Single Peer, and Multi-Peer condition), our sample exceeds this threshold. 

We recruited our participants via the Prolific crowdsourcing platform known to provide high quality crowdsourced data \cite{Douglas2023DataSONA}. Using Prolific's pre-screening tools, we targeted higher-education students and academic professionals. This demographic was selected because university students and early-career researchers represent a population that frequently experiences high levels of academic pressure, professional burnout, and imposter syndrome. Making them a valid target group for our context-specific prompts regarding workplace and academic stress and allowing us to safely assess the effectiveness of social modelling in lowering everyday professional stigma without targeting vulnerable clinical populations.

\subsection{The Interaction Task}
Participants will be randomly assigned to one of three agent configurations (Baseline, Single Peer, or Multi-Peer). Across all conditions, the interaction will last approximately 10 to 15 minutes and is structured around the three conversational stages. To establish a baseline of psychological safety in the multi-agent conditions, the personas are designed to maintain a controlled and non-judgmental conversational tone. We utilise these specific traits as a controlled characteristic to ensure the social modelling remains consistently supportive throughout the session. In the following section, we detail how the interaction task is implemented by our prototype, and we specify the behaviours of the chatbots.

\subsection{The Prototype}

\subsubsection{The System}
We implemented the prototype as a browser-based experimental system composed of a React/TypeScript frontend and a Go backend. The interface guided participants through onboarding and consent, a pre-interaction questionnaire, the conversational task, a post-interaction questionnaire, and a completion screen. The backend managed session state, randomly assigned participants to one of the three study conditions, and logged transcripts and questionnaire responses for later analysis.

The prototype was designed as a tightly controlled experimental instrument rather than a general-purpose support tool. Across conditions, participants experienced the same three-stage interaction structure and the same overall interface; the only systematic manipulation was the presence and number of peer agents. To preserve ecological validity without undermining the experimental manipulation, the interface introduced subtle typing delays before bot turns and brief comfort check-ins after each conversational stage.

\begin{figure}[t]
    \centering
    \fbox{\parbox[c][3.1cm][c]{0.9\columnwidth}{\centering \textit{Placeholder screenshot: onboarding and consent flow}}}
    \caption{Placeholder screenshot of the onboarding and consent flow.}
    \Description{Screenshot placeholder for the participant onboarding experience.}
    \label{fig:prototype-onboarding}
\end{figure}

\begin{figure}[t]
    \centering
    \fbox{\parbox[c][3.1cm][c]{0.9\columnwidth}{\centering \textit{Placeholder screenshot: chat interface with host and peer agents}}}
    \caption{Placeholder screenshot of the multi-agent chat interface used during the interaction.}
    \Description{Screenshot placeholder for the experimental chat interface.}
    \label{fig:prototype-chat}
\end{figure}

\subsubsection{The Chatbots}
The interaction relied on two distinct agent roles. The host facilitator, Vieno, was implemented as a live LLM-powered agent using Anthropic Claude Sonnet 5 via OpenRouter. To keep Vieno aligned with the study's social modelling goal, the system prompt constrained the model to brief, non-evaluative acknowledgements that reflected the participant's most recent message. The host was therefore responsible for maintaining conversational flow rather than for interpreting, advising, or introducing new topics. If generation failed or timed out, the system fell back to a fixed acknowledgement (``Thanks for sharing that.'') so that the session could continue without interruption.

The peer agents, Sam and Charlie, were not generated live. Instead, they were scripted deterministically from condition-specific dialogue scripts so that their content and timing remained constant across participants. In the single-peer condition, Sam contributed a short neutral opener and then progressively more personal disclosures. In the multi-peer condition, both Sam and Charlie contributed comparable turns, creating a group-like conversational setting in which participants observed peer vulnerability before being invited to respond. To preserve the illusion of a live conversation, the interface introduced word-count-based typing delays before bot turns, with a maximum delay of approximately 10 seconds.

The conversation itself was structured around three escalating stages. In the early stage, participants were asked about everyday professional or academic challenges. In the mid stage, the conversation moved toward a recent experience of feeling overwhelmed while not being able to share it. In the late stage, participants were invited to consider a more personal professional or academic struggle. In the baseline condition, the participant answered the host directly. In the peer conditions, the scripted peers disclosed progressively more personal experiences before the host invited the participant to respond, giving the participant an observable social model of disclosure before their own turn.

\section{Data Analysis}

To evaluate the textual data generated during the interaction, we will utilise a deductive qualitative coding approach. The chat transcripts will be analyzed by two independent researchers using a strict qualitative codebook to ensure objective measurement of state disclosure. Inter-rater reliability will be calculated using Cohen's Kappa. The codebook tracks three primary behavioral dimensions:

\textbf{Depth of Vulnerability:} Measures the thematic content of the disclosure, scaling from superficial complaints (\textit{VUL-LOW}, e.g., external logistics) to severe stigma markers (\textit{VUL-HIGH}, e.g., admissions of perceived incompetence or imposter syndrome).

\textbf{Linguistic Markers:} Measures structural masking behaviors. We specifically track instances of masking (\textit{LIN-GUARD}, e.g., third-person distancing) versus nervous disclosure (\textit{LIN-HESIT}, e.g., trailing sentences and self-minimizing language).

\textbf{Social Modeling Efficacy:} To test our between-subjects hypotheses, we track whether users actively utilise the peer agents' vulnerability as a social script. Codes capture both validation (\textit{SOC-ECHO}, where the user explicitly references a bot's narrative) and resistance (\textit{SOC-DEFLECT}, where the user ignores the modeled vulnerability).

\section{Discussion}

This research contributes a novel design framework for mental health interactions by shifting away from monolithic dyadic agents toward multi-agent peer personas. We expect our results to demonstrate that social modelling is not merely a psychological theory, but a functional design requirement for eliciting honest, stigmatised disclosures in HCI. By reducing the evaluative anxiety inherent in monolithic chatbots, we pave the way for more human-centric, trustworthy health informatics.

In exploring how peer disclosure shapes participant responses (\textbf{RQ1}), our premise rests on the dual mechanisms of social proof and anonymity dilution. By witnessing a group of agents engage in vulnerability, the user receives a stronger signal that such behaviour is the ``norm'' in this specific digital space \cite{Lugrin2021HandbookVol1, Lugrin2022HandbookVol2}. As the conversation shifts to a group dynamic rather than a dyadic interrogation, the user may feel less like a clinical subject under a microscope and more like a participant in a communal dialogue.

Furthermore, our investigation into how disclosure unfolds temporally (\textbf{RQ2}) is grounded in the principles of progressive engagement. As demonstrated by Song et al., trust and empathetic connection in human-AI cooperation are not instantaneous but accrue over the course of an interaction \cite{Song2025Cooperation}. Similarly, Sien and McGrenere note that digital storytelling platforms require support for gradual engagement to help users overcome their initial reservations \cite{Sien2025DigitalHumanLibrary}. By tracking users stage-by-stage, we expect to empirically capture this trajectory, showing how sustained social modelling dismantles the user's initial linguistic masking.

Finally, regarding the mitigation of evaluative anxiety (\textbf{RQ3}), we anticipate that observing peer vulnerability will alleviate the pressure of direct assessment. Traditional dyadic agents often unintentionally mimic clinical interviews, which can elevate a user's apprehension of being judged. By contrasting real-time comfort trajectories with post-interaction metrics, we aim to demonstrate that a communal, multi-agent environment acts as a psychologically safe buffer, mitigating the fear of negative evaluation even when users are prompted to share deeply personal struggles.

\section{Conclusion}
In this paper, we presented a 3x3 mixed-methods study investigating the impact of multi-agent social modelling on user disclosure over time. By shifting the user's role from a participant being questioned to an observer of a safe social dialogue, this design framework lowers the barriers created by self-stigma and mitigates the evaluative anxiety of monolithic interactions.

\begin{acks}

\end{acks}

\appendix
\section{LLM Prompt(s) Used in the Prototype}
\label{appendix:prompts}

This appendix contains the prompt used to constrain the live host facilitator in the prototype. The prompt was designed to keep the agent in a brief acknowledgement role and to prevent it from offering advice, interpretation, or new conversational topics.

\subsection{Host Acknowledgement Prompt}
\label{appendix:systemprompt}

\begin{quote}
You are Vieno, hosting a short peer support chat.\newline
A participant has just written a message. Reply with a brief acknowledgement.\newline
\newline
Rules:\newline
- Write exactly one sentence, at most 25 words.\newline
- Reflect only what the participant actually wrote. Never mention anything they did not say.\newline
- Do not give advice, suggestions, or next steps.\newline
- Do not evaluate, praise, diagnose, or interpret what they said.\newline
- Do not introduce a new topic and do not ask a question.\newline
- Do not include internal or system XML tags in your response.\newline
- Write plainly, as a person would in a chat.\newline
\end{quote}

\subsection{Notes on Prompt Design}
The purpose of this prompt was not to make the host conversationally rich, but to preserve the integrity of the experimental manipulation. By constraining the LLM to short acknowledgements, the prototype kept the participant's experience focused on the presence or absence of peer disclosure rather than on the host's own interpretive behaviour.


\bibliographystyle{ACM-Reference-Format}
\bibliography{references}

\end{document}
